\documentclass[10pt,twocolumn]{article}

\usepackage[spanish, activeacute]{babel}
\usepackage[utf8]{inputenc}
%\usepackage{oz2tex}
\usepackage{graphicx}
\usepackage{blindtext}
\usepackage{amsmath}
\usepackage{amsfonts}
\usepackage{float}
\usepackage{url}
\usepackage{anysize} %Para usar márgenes
\usepackage{hyperref}
\marginsize{1cm}{1cm}{1.5cm}{1.5cm} %{izquierda}{derecha}{arriba}{abajo}.
\usepackage{color}
\usepackage{colortbl}
\usepackage{moreverb}

\usepackage{array}
\definecolor{naranja}{rgb}{1,0.5,0} % valores de las componentes roja, verde y azul (RGB)
\definecolor{rojo}{rgb}{1,0,0}
\definecolor{SteelBlue}{rgb}{0.3,0.5,0.7}
\definecolor{violet}{rgb}{0.35,0,0.34}

\usepackage{listings}
\usepackage{listingsutf8}
\usepackage{amsmath, amsthm, amssymb}
\lstset{ %
  language=python,
  basicstyle=\small,           % the size of the fonts that are used for the code
  numbers=none,
  numberstyle=\footnotesize,          % the size of the fonts that are used for the line-numbers
  numbersep=4pt,                  % how far the line-numbers are from the code
  backgroundcolor=\color{white},      % choose the background color. You must add \usepackage{color}
  breaklines=true,                % sets automatic line breaking
  breakatwhitespace=false,        % sets if automatic breaks should only happen at whitespace
  extendedchars=true,
  inputencoding=utf8, 
    tabsize=2,
    aboveskip=5pt,
    upquote=true,
    showstringspaces=false,
    frame=single,
    showtabs=false,
    showspaces=false,
    showstringspaces=false,
    keywordstyle=\color{blue}\ttfamily\bfseries,
    stringstyle=\color{red}\ttfamily,
    commentstyle=\color{violet}\ttfamily,
    morecomment=[l][\color{magenta}]{\#}, 
    literate={á}{{\'a}}1 {é}{{\'e}}1 {í}{{\'i}}1 {ó}{{\'o}}1 {ú}{{\'u}}1 {'}{{'}}1 {ñ}{{\~n}}1 {&}{{\&}}1 ,                                        
}

\lstset{emph={%  
    define-struct, in, then,else, letrec, symbol=,%
    },emphstyle={\color{blue}\ttfamily\bfseries}%
}%

\usepackage{tikz}
\usepackage{diagbox}

\usepackage{fancyhdr} % Para personalizar encabezados y pies de página
\usepackage{lipsum} % Para generar texto de muestra
\fancyfoot[C]{Página \thepage} % Número de página en el centro del pie de página
\renewcommand{\headrulewidth}{0pt} % Eliminar la línea del encabezado

% Aplicar el estilo de encabezado y pie de página
\pagestyle{fancy} % Para personalizar encabezados y pies de página
\usepackage{lipsum} % Para generar texto de muestra

% Definir el encabezado y pie de página
\fancyhead{} % Limpiar todos los campos de encabezado
\fancyhead[C]{
  \begin{tabular}{|c|c|c|c|}
    \hline
      \begin{tikzpicture}[baseline=(current bounding box.center)]
        \draw (0,0) circle (0.25cm);
      \end{tikzpicture} &
      \begin{tikzpicture}[baseline=(current bounding box.center)]
        \fill[black] (0,0) circle (0.25cm);
      \end{tikzpicture} &
      \begin{tikzpicture}[baseline=(current bounding box.center)]
        \draw (0,0) circle (0.25cm);
      \end{tikzpicture} &
      \begin{tikzpicture}[baseline=(current bounding box.center)]
        \draw (0,0) circle (0.25cm);
      \end{tikzpicture} \\
      \hline
  \end{tabular}} % Contenido del encabezado en la izquierda

% Aplicar el estilo de encabezado y pie de página

\title{\includegraphics[scale=0.2]{../imagenes/univalle.png} \\ Fundamentos de compilación e interpretación de lenguajes de programación \\ Primer exámen parcial, Duración 2 horas \vspace{-0.5cm}} 
\author{Carlos Andres Delgado S, Msc \\ \texttt{ carlos.andres.delgado@correounivalle.edu.co}}
\date{\vspace{-0.2cm}03 de Abril de 2025}


\newcommand{\raya}{\underline{\hspace{4cm}}}

\begin{document}

\maketitle


\thispagestyle{fancy}


\begin{enumerate}
\item {
  ¿Que se debe tener en cuenta al intentar implementar procedimientos recursivos con let?
  \begin{enumerate}
      \item Se requiere un nuevo tipo de expresión, ya que podríamos tener definiciones recursivas de ligaduras %correcto
      \item let sólo sirve para ligar valores que no sean procedimientos
      \item Se debe cambiar la forma de construir de clausura para tener el ambiente actual y no el anterior
      \item Al extender del ambiente anterior, los procedimientos no se conocen a sí mismos %correcto
      \item Las clausuras no contienen el cuerpo del procedimiento, por lo cual no puede conocerlo
  \end{enumerate}
}

\item {
	En un TAD las funciones que contamos son constructores, predicados y extractores que permiten:

	\begin{enumerate}
		\item Extractores: Permiten conocer que variante pertenece un tipo de dato
		\item Los predicados y extractores se denominan procedimientos observadores %correcto
		\item Constructor: Permite crear un elemento perteneciente al TAD %correcto
		\item Predicado: Permite insertar un elemento nuevo al TAD
		\item Un programador puede usar los predicados para conocer los valores internos de un TAD
	\end{enumerate}
}


\item {
Indique cuales de las siguientes opciones son correctas con respecto a la receta de diseño de tipos de datos recursivos
	\begin{enumerate}
		\item Incluir una función de validación de datos, esto para asegurar que los constructores trabajan apropiadamente 
		\item Incluir un extractor para cada pieza de dato %correcta
		\item Incluir un constructor para cada clase de dato %correcta
		\item Incluir una representación de listas o procedimientos para manejar el dato en Racket fácilmente
		\item Incluir un predicado para cada clase de dato en el tipo de dato %correcta
  \end{enumerate}
}

\item {  
¿Que es sintaxis concreta o representación externa?
\begin{enumerate}

\item Es el código binario generado por un compilador de un lenguaje.
\item Es código en un lenguaje de bajo nivel.
\item Es código en un lenguaje de alto nivel.
\item Es la representación de un tipo de dato a partir de su gramática.
\item Es una estrategia para transformar código en un lenguaje de bajo nivel a uno de alto nivel.
  \end{enumerate}
}

\item {
Un parser es un procedimiento que:
	\begin{enumerate}
		\item Determina si una variable esta ligada o no
		\item Permite buscar los valores de ligaduras en los ambientes
		\item Transforma sintaxis abstracta a concreta
		\item Extrae las unidades significativas de la sintaxis concreta
		\item Transforma sintaxis concreta a sintaxis abstracta %correcto
	\end{enumerate}
}
  

\item {
 Dada la siguiente definición de TAD
  \begin{lstlisting}
  (define-datatype dato dato?
    (datoA (x number?))
    (datoB (y symbol?) (z (list-of symbol?) (w dato?)))
  )
  \end{lstlisting}
  ¿Que declaraciones son correctas?
	\begin{enumerate}
	  \item (datoB 'x '(a b c) (datoA 5)) %Correcto
    \item (datoA 'x)
    \item (datoB 'z '(t u v) (datoB 'y '(a b c) (datoA 3))) %Correcto
    \item (datoB 'p (datoB 'x '(a b c) (datoA 3)))
    \item (datoB 'z '(a b c))
  \end{enumerate}
}

\item {
  
  ¿Que limitación tiene SLLGEN con respecto a la gramáticas?
  \begin{enumerate}

\item No se puede tener variables mutables
\item Sólo puede reconocer números enteros como datos
\item No se puede diferenciar casos que inician igual por la izquierda, ejemplo var x  con var x = 10 %correcto
\item Se deben hacer todos los datatypes con define-datatype para que el interprete funcione
\item No se pueden tener expresiones de más de una línea
  \end{enumerate}

}

\item {
Indique que es cierto con respecto a un tipo abstracto de dato (TAD):
	\begin{enumerate}
		\item Si la implementación interna cambia, el programador debe ajustar sus funciones
		\item Tienen implementación e interfaz %correcto
		\item Si la implementación interna cambia, el programador no nota la diferencia %correcto
		\item Tienen Representación interna y predicados
		\item Tienen dos casos, TAD vacío que es el base y TAD extendido que permite extender a partir de una definición anterior
	\end{enumerate}

}

\item {
Cuando cambiamos la representación interna de un tipo de dato (listas, procedimientos, bits, etc) ¿Que sucede con las funciones externas que usan la interfaz?

	\begin{enumerate}
		\item Cambian, es necesario reestructurar las funciones
		\item Cambian levemente, es necesario ajustar los condicionales de las funciones
		\item No cambian, pero el cambio de la representación puede implicar nuevas funciones auxiliares
		\item No cambian, el cambio de representación no las afecta %correcto
		\item Cambian, pero solo es necesario cambiar los condicionales y los llamados a los extractores
	\end{enumerate}
}

\item {
Un ambiente es un tipo de dato recursivo que:
	\begin{enumerate}
		\item El caso de ambiente vacío tiene una lista de identificadores y una lista de valores
		\item Permite representar sintaxis abstracta desde sintaxis concreta
		\item Tiene dos casos ambiente vacío y ambiente extendido %correcta
		\item Tiene tres casos: ambiente vacío, ambiente extendido y ambiente sobreextendido
		\item El caso de ambiente extendido contiene lista de identificadores, lista de valores y ambiente que extiende %correcta
	\end{enumerate}
}

\end{enumerate}
Conteste las preguntas de 11 a 15 de acuerdo al siguiente código:

Considere el ambiente inicial env0 \texttt{(x,y,z,f), (4,2,1, closure '(a,b) +(2,+(a,b)) empty-env)} Dibuje los ambientes para la expresión:
\begin{lstlisting}
let
 a = let x = +(x,z) in (f x y)
 b = let x = +(x,y) in (f x z)
 c = let x = +(x,x) in (f y z)
 in
 letrec
  g(x,y) = if >(y,0) then +(*(2, x),(g +(x,1) -(y,1))) else let a = +(a,x) in (f a b)
  in
    ( proc (n,m) (g n m) b c)
\end{lstlisting}

\begin{enumerate}
\setcounter{enumi}{10}
\item {
¿Cuantos ambientes del procedimiento f se generan?
\begin{enumerate}
	\item 2
	\item 3
	\item 4 %correcto
	\item 5
	\item 6
\end{enumerate}
}

\item {
Para los valores de a, b, c  es cierto que:
\begin{enumerate}
	\item b = 6
	\item a = 3 
	\item b = 5 %correcto
	\item a = 6
	\item a = 7
\end{enumerate}
}

\item {
El valor final de la expresión es:
\begin{enumerate}
	\item 159
	\item 144 %correcto
	\item 129
	\item 122
	\item 153
\end{enumerate}
}

\item {
 a =  let x = +(x,z) in .... ¿Donde extiende el ambiente que contiene x = +(x,y)?
\begin{enumerate}
	\item Ambiente del procedimiento de f que se invoca
	\item Ambiente inicial env0 %correcto
	\item Ambiente vacio
	\item Ambiente que contiene a,b,c
	\item Ambiente recursivo de g(x,y)
\end{enumerate}
}

\item {
Para los valores de a, b, c  es cierto que:
\begin{enumerate}
	\item c = 6 %correcto
	\item a = 6
	\item c = 5 
	\item a = 9 %correcto
	\item a = 7
\end{enumerate}
}

\end{enumerate}
\end{document}
